\section{Thesis Statement}
\label{sec:intro:statement}

\textit{By explicitly incorporating dynamic constraints into the motion generation process---through sampling-based kinodynamic planning, learned generative models, and dynamics-informed trajectory seeding---robots can perform a substantially broader range of manipulation tasks while operating closer to their true physical limits.}

This thesis advances robot manipulation by developing a progression of methods that expand the scope of manipulation beyond geometric pick-and-place to include dynamic non-prehensile skills and motions governed by object inertial properties, impulsive contact, and motion-induced forces.
Across all three research themes, dynamic constraints are treated not as a complexity to be avoided but as structured information to be modeled, learned, and exploited.
